\section{Learning Objectives}

After completing this chapter, readers will be able to:
\begin{itemize}
    \item Apply software engineering principles to ML code
    \item Structure ML projects effectively
    \item Implement modular and testable ML code
    \item Use design patterns for ML systems
    \item Create reproducible training pipelines
\end{itemize}

\section{Code Organization for ML Projects}

\subsection{Project Structure}

\begin{lstlisting}[style=shell]
ml_project/
├── data/
│   ├── raw/
│   ├── processed/
│   └── features/
├── models/
├── notebooks/
├── src/
│   ├── data/
│   │   ├── make_dataset.py
│   │   └── preprocessing.py
│   ├── features/
│   │   └── build_features.py
│   ├── models/
│   │   ├── train.py
│   │   ├── predict.py
│   │   └── evaluate.py
│   └── utils/
├── tests/
├── config/
├── requirements.txt
└── README.md
\end{lstlisting}

\section{Model Development Patterns}

% Placeholder content

\section{Testing ML Code}

% Placeholder content

\section{Chapter Summary}

% Placeholder summary

\section{Exercises}

\begin{exercise}
Refactor a Jupyter notebook into a structured, testable codebase.
\end{exercise}
